% quiz-macros.tex -- shared package setup for quiz-template.tex (exam class).
%
% NOTE ON PATHS: every template in this repository is stored one level up
% from where it is compiled (see README.md). All relative \input paths
% below are resolved against the folder the document is compiled in,
% which sits one directory above this file.

% bm: bold vectors (\bm{v}) for physics-style notation in quiz math.
\usepackage{bm}

% --- exam-class patch: point labels next to each \item ---------------------------
% Redefines exam.cls's internal \item@points@pageinfo so the "N pts" label
% is typeset in a group and inserted into the \@labels box -- the box
% exam.cls uses (since v2.218) to attach page-info to the right \item even
% when a question starts with a list/verbatim environment (those
% environments clobber \everypar, which is why the old approach broke).
% \makeatletter is required so the @-names below resolve; the matching
% \makeatother below restores normal catcodes for the rest of the document.
\makeatletter
\ExplSyntaxOn
\def\item@points@pageinfo{%
  \item 
  \if@placepoints
    \if@bonus
      \def\padded@point@block{%
        \begingroup
          \@placepointstrue
          \bonuspoint@block
        \endgroup
      }%
    \else
      \def\padded@point@block{%
        \begingroup
          \@placepointstrue
          \point@block
        \endgroup
      }%
    \fi
    \if@pointsdropped
      % Do nothing!
    \else
      \if@bonus
        \if@bonusqformat
          \ifx\ques@ref\@queslevel
            % Do nothing
          \else
            \setup@point@toks
          \fi
        \else
          \setup@point@toks
        \fi
      \else
        \if@qformat
          \ifx\ques@ref\@queslevel
            % Do nothing
          \else
            \setup@point@toks
          \fi
        \else
          \setup@point@toks
        \fi
      \fi
    \fi
    \global \@placepointsfalse
  \fi
  % Version 2.218$\beta$, 2007-10-31 changes:
  % Instead of appending
  %   \the \pageinfo@commands \the \point@toks
  % to \everypar, we insert them into the box \@labels.  This corrects
  % the problem that arose when a question (or part, etc.) begins with
  % a list environment (including verbatim, flushleft, center,
  % flushright, and possibly others that are implemented as trivlist
  % environments).  The \item command in those environments throws
  % away the previous contents of \everypar, and so the tokens \the
  % \pageinfo@commands \the \point@toks didn't get inserted where we
  % expected.  List environments *do* preserve the contents of the box
  % \@labels, though.
  \global\setbox\g__block_labels_box\hbox{\unhbox\g__block_labels_box %<----------use new name ...
    \the \pageinfo@commands
    \the \point@toks}%
}
\ExplSyntaxOff
\makeatother

% --- Tagged PDF: table headers -----------------------------------------------------
% Tell tagpdf to mark row 1 of every table as a header row, so screen
% readers announce it as a table header.
\tagpdfsetup{table/header-rows={1}}

% --- Name / Student ID box --------------------------------------------------------------
% Prints the answer instructions plus fill-in lines for Name and Student
% ID at the top of the quiz sheet.
\newcommand{\NameSection}{\vspace{-2in}\begin{center}
\fbox{\fbox{\parbox{5.5in}{\centering
Answer the questions in the spaces provided on the
question sheets. If you run out of room for an answer,
continue on the back of the page.}}}
\end{center}
\vspace{0.2in}
\makebox[\textwidth]{Name: \enspace\hrulefill}
\makebox[\textwidth]{Student ID: \enspace\hrulefill}}

% --- Points summary ----------------------------------------------------------------------
% "This quiz is worth \numpoints points ... 15 minutes" plus the exam-
% class \pointtable listing the points per question.
\newcommand{\PointsSection}{
  This quiz is worth \numpoints\ points and consists of \numquestions\ questions. You have 15 minutes to complete the quiz.

  \begin{center}
    \pointtable[v][questions]
  \end{center}
}

% --- Math -------------------------------------------------------------------------------
\usepackage{amsmath,amssymb}

% --- Shared course macros ---------------------------------------------------------------
% Root macros.tex: math shorthands used across all course documents
% (\RR, \NN, \ZZ, \QQ, \grad, \curl, \tr, \der, \pder, ...).
\newcommand{\RR}{\mathbb{R}}
\newcommand{\NN}{\mathbb{N}}
\newcommand{\ZZ}{\mathbb{Z}}
\newcommand{\QQ}{\mathbb{Q}}
\newcommand{\vv}{\vec{v}}
\newcommand{\ww}{\vec{w}}
\newcommand{\xx}{\vec{x}}
\newcommand{\yy}{\vec{y}}
\newcommand{\zz}{\vec{z}}
\newcommand{\uu}{\vec{u}}
\newcommand{\basis}[1]{\vec{e_{#1}}}
\newcommand{\id}{\mathrm{id}}
\newcommand{\CC}{\mathbb{C}}
\renewcommand{\div}{\operatorname{div}}
\newcommand{\grad}{\operatorname{grad}}
\newcommand{\curl}{\operatorname{curl}}
\newcommand{\tr}{\operatorname{tr}}
\newcommand{\der}[2]{\dfrac{d #1}{d #2}}
\newcommand{\pder}[3][]{\dfrac{\partial^{#1} #2}{\partial #3}}

